% =============================================================================
% Wprowadzenie. Plik edytowany ręcznie -- skrypt analiza_odpowiedzi_naukowa.py go NIE nadpisuje.
% Liczby: definicje.tex; tabele: \TabDir (CSV); rysunki: \FigDir; literatura: references.bib.
% =============================================================================
\section{Wprowadzenie}

Przypominanie sobie informacji z pamięci wzmacnia ich późniejsze utrwalenie silniej niż ponowne
ich przeczytanie (efekt testowania), a nauka rozłożona w czasie jest zapamiętywana lepiej niż nauka
skumulowana (efekt rozłożenia w czasie) \citep{roediger2006,cepeda2006,cepeda2008}. W randomizowanych
badaniach z udziałem lekarzy w trakcie kształcenia tzw. edukacja rozłożona w czasie (ang.
\emph{spaced education}), w której krótkie pytania wielokrotnego wyboru przesyła się cyklicznie
przez tygodnie lub miesiące, poprawiała przyswajanie i utrwalanie wiedzy
\citep{kerfoot2007urology,kerfoot2007mededuc,kerfoot2009}.

W szkoleniach organizacyjnych stosuje się \emph{mikrouczenie} (ang. \emph{microlearning}), czyli
krótkie jednostki szkoleniowe, często w formie quizów. Dowody jego skuteczności pochodzą jednak głównie
z badań na małych próbach lub z tzw. szarej literatury \citep{taylor2022}, a ewaluacje obowiązkowych
szkoleń organizacyjnych, np.\ z zakresu świadomości bezpieczeństwa, często wykazują efekty niewielkie
albo zanikające z upływem czasu \citep{reinheimer2020,lain2022,ho2025}. Otwarte pozostaje więc
pytanie, czy wiedza sprawdzana w trakcie takich szkoleń jest nabywana i utrwalana.

Skuteczność szkoleń organizacyjnych ocenia się zwykle przez transfer, rozumiany jako zastosowanie
nabytej wiedzy i umiejętności w pracy \citep{baldwin1988,blume2010}. Warunkiem transferu jest jednak
wcześniejsze nabycie i utrwalenie wiedzy, a ten etap można zweryfikować na podstawie danych
gromadzonych przez platformy mikrouczenia. Platformy te zapisują każdą odpowiedź wraz z czasem
i wynikiem, a to samo pytanie wraca do tego samego pracownika wielokrotnie. Kolejne próby tworzą
zatem serię pomiarów tej samej wiedzy u tej samej osoby, co pozwala ocenić, czy poprawność rośnie,
czy wiedza utrzymuje się między powtórkami i czy błędy się utrwalają. Dane na poziomie pojedynczych
odpowiedzi wykorzystuje się w podobny sposób w modelach ucznia rozwijanych dla inteligentnych systemów
tutorskich i aplikacji do nauki języków \citep{corbett1994,pavlik2009,settles2016,tabibian2019,koedinger2023}.
W niniejszym badaniu podejście to zastosowano do walidacji wiedzy nabywanej w ramach mikrouczenia
w organizacji.

Badanie przeprowadzono w banku, w którym pracownicy odpowiadali na krótkie pytania wielokrotnego
wyboru dotyczące wymogów regulacyjnych, produktów bankowych, umiejętności miękkich oraz, w przypadku
kadry kierowniczej, zarządzania, wyświetlane przez aplikację na komputerze w trakcie dnia pracy.
Pytania wracały do pracownika w stałym, nieadaptacyjnym odstępie, a po każdej odpowiedzi wyświetlano
odpowiedź poprawną. Na podstawie analizy odpowiedzi z okresu trzech lat,  sformułowano trzy pytania badawcze:
(PB1) czy odsetek poprawnych odpowiedzi rośnie wraz z kolejnymi przypomnieniami tego samego pytania
przez tego samego pracownika;
(PB2) jak poprawność zależy od wyniku poprzedniej próby i od czasu, który od niej upłynął;
(PB3) czy powtórzone błędy są trwałe, co wskazywałoby raczej na utrwalone błędne przekonania niż na
zgadywanie, mimo informacji zwrotnej po każdej odpowiedzi.
